\paragraph{A.5c (Refuted: arbitrary ramified base change under the module-category definition).}
The universal statement at 15:36 is false if ``Deuring category'' retains
the module-category conditions (b),(c) of 15:16 after an arbitrary
commutative base change. This remains false when the category over the
new base is the actual base change of a global maximal-order category,
with the base-changed positive determinant object. Its category of
determinant-rigidified autoequivalences need not be a Picard category.
We give a counterexample over $R=\mathbb F_9$ in which even the group
of isomorphism classes is noncommutative.

\begin{proof}
We take the maximal order $O\subset B_3$ of A.1a and let $C$ be its
category of right modules locally free of rank one over
$\operatorname{Spec}\mathbb Z$. Put
\[
 R=\mathbb F_3[\alpha]/(\alpha^2+1),\qquad B=O\otimes_{\mathbb Z}R.
\]
Since $2$ is invertible at $3$, the localization of $O$ has generators
$i,j$ with $i^2=-1$, $j^2=-3$, and $ij=-ji$. Consequently $B$ has
orthogonal idempotents
\[
 e_0=(1+i/\alpha)/2,\qquad e_1=(1-i/\alpha)/2,
 \qquad x=e_0je_1,\quad y=e_1je_0.
\]
These form an $R$-basis. Their multiplication is specified by
$e_0+e_1=1$, $x=e_0xe_1$, $y=e_1ye_0$, and all products of two
elements of $Rx+Ry$ being zero. In coordinates the product and the
base-changed reduced norm are
\begin{align*}
 (a,b,c,d)(a',b',c',d')
   &=(aa',bb',ac'+cb',bd'+da'),\\
 N(ae_0+be_1+cx+dy)&=ab.
\end{align*}
The displayed polynomial is the reduction of the quaternion reduced
norm, not the determinant of the four-dimensional regular
representation (which would be its square).
The centre of $B$ is exactly $R$: commuting with $e_0$ forces $c=d=0$,
and commuting with $x$ then forces $a=b$. In particular no central
radical elements are being omitted.

This is the actual base change of $C$. Indeed, for finite projective
$O$-modules $P,Q$, base change identifies
\[
 \operatorname{Hom}_O(P,Q)\otimes R
 \simeq \operatorname{Hom}_B(P\otimes R,Q\otimes R).
\]
Every $P$ is free of rank one locally at $3$, so $P\otimes R$ is a
free rank-one $B$-module. Conversely $O\otimes R=B$ occurs, and over
the field $R$ every module locally free of rank one over
$\operatorname{Spec}R$ is free. Thus $C_R$ is equivalent to the
one-object $R$-linear category with endomorphism algebra $B$, and to
the full category of rank-one $B$-modules. This verifies conditions
(b),(c) of 15:16 without replacing $C_R$ by an unrelated category.

We specify its determinant rigidification carefully. For a free
rank-one right $B$-module $P$, its norm line is the contracted product
of the torsor of bases of $P$ with $R$ under $N:B^*\to R^*$.
Changing a basis by $b\in B^*$ acts on this line by $N(b)$.
This construction is the base change of the reduced-norm-line
construction for $O$. Choose the object $X=O$ before base change;
its identity basis has norm one and identifies its determinant object
with the positive reference object $T$ of A.6a. It therefore identifies
$\nu(X_R)$ with $T_R$, with no additional scalar choice.

On the one-object model an $R$-algebra automorphism $\rho$ of $B$
preserving the polynomial $N$ gives an admissible autoequivalence
$F_\rho$. Equivalently it gives the invertible bimodule with one
action twisted by $\rho$, free of rank one on both sides.
Since $N\circ\rho=N$, its induced functor on the determinant object
is the identity in the preceding coordinates. We rigidify it by
$u_\rho=\operatorname{id}_{T_R}$. These are specified rigidifications,
not merely an assertion that its transported norm line happens to
be trivial. Products have the same identity rigidification.
A natural isomorphism $F_\rho\to F_\kappa$ is a unit $b\in B^*$
satisfying $b\rho(a)=\kappa(a)b$ for all $a\in B$; its map on
the determinant line is multiplication by $N(b)$. It respects these
chosen rigidifications if and only if $N(b)=1$.

Set $t=1+\alpha$. Its square is $-\alpha$, its fourth power is $-1$,
and its order in $R^*$ is eight; hence $t$ is a nonsquare. Define
two $R$-algebra automorphisms by
\[
 S(e_0)=e_1,\quad S(e_1)=e_0,\quad S(x)=y,\quad S(y)=x,
 \qquad
 T_t(e_i)=e_i,\quad T_t(x)=tx,\quad T_t(y)=y.
\]
The multiplication table verifies these assertions directly, and
both preserve $N$. Write $\rho=S\circ T_t$, $\kappa=T_t\circ S$.
For the unit $d=te_0+e_1$ we have
\[
 \kappa=\operatorname{Int}(d)\circ\rho,\qquad N(d)=t.
\]
If $b\rho(a)=\kappa(a)b$ for every $a$, then $d^{-1}b$ commutes
with $\rho(B)=B$. The centre calculation gives $b=cd$ for
$c\in R^*$, and therefore
\[
 N(b)=c^2t\ne1.
\]
Thus the unrigidified products have an isomorphism, but no isomorphism
preserves their specified determinant rigidifications. The classes
of $(F_S,1)$ and $(F_{T_t},1)$ do not commute. A Picard category
necessarily has commutative $\pi_0$, so the proposed Picard structure
and hence the assertion of a canonical equivalence \emph{as Picard
categories} fail for this base change. Taking
$(C',T')=(C_R,T_R)$ already supplies the counterexample.

The script \texttt{verify\_ramified\_picard\_counterexample.py} checks
the multiplication, the automorphisms and the centre, enumerates all
$9^4=6561$ elements and all $8^2\cdot9^2=5184$ units, and finds
precisely the eight intertwiners $cd$. Their norms are the four
nonsquares of $R^*$; none is one. The argument above proves the result
independently of this finite verification.
\end{proof}

\paragraph{A.5d (Proved qualification of the counterexample).}
The category in A.5c satisfies the broad module-category definition
of 15:16. The algebra $B$ is not Azumaya: its nonzero radical
$Rx+Ry$ has square zero. It therefore does not satisfy the separately
bracketed maximal-order condition (a) over this field. If the phrase
``Deuring category'' is restricted to strict maximal-order categories
or to Azumaya categories after base change, A.5c is outside that
restricted domain and A.5a applies there. The universal quantifier
over arbitrary commutative base rings can thus be retained only
with an explicit admissibility condition on the resulting category;
under conditions (b),(c) alone it is refuted, rather than unresolved.
